% Part: normal-modal-logic
% Chapter: axioms-systems
% Section: normal-logics

\documentclass[../../../include/open-logic-section]{subfiles}

\begin{document}

\olfileid{nml}{prf}{nor}

\olsection{نارمل امکانی منطقاں}

امکانی !!{formula}s دے کجھ سیٹ ایسے نیں جیہناں نوں کسے مسلّمیاں دے سیٹ توں اخذ ہون والے !!{formula}s وجوں پچھاننا سوکھا نہیں۔ اسیں چاہندے آں کہ امکانی منطقاں کجھ خاص قاعدے پالن۔ سبھ توں پہلاں، جو کجھ اسیں کلاسیکی قضیاتی منطق وچ !!{derive} کر سکدے آں، اوہ اجے وی !!{derivable} ہونا چاہیدا اے۔ ایہ خیال رکھنا اے کہ ہُن !!{formula}s وچ
\iftag{prvBox}{$\Box$\iftag{prvDiamond}{
    تے~}{}}{}\iftag{prvDiamond}{$\Diamond$}{} وی آ سکدے نیں۔ ایس لئی اسیں شرط لاؤندے آں کہ امکانی منطق وچ ہمیشہ صادق فارمولیاں توں بدل نال بنیاں ساریاں مثالاں ہون، تے ایہ قاعدۂ اثباتِ مقدم تحت بند ہووے۔

\begin{defn}
  \emph{امکانی منطق} امکانی !!{formula}s دا ایسا سیٹ~$\Sigma$ اے جیہڑا
  \begin{enumerate}
  \item سارے ہمیشہ صادق فارمولے رکھدا اے؛
  \item بدل تحت بند اے، یعنی جے $!A \in \Sigma$ تے
    $!D_1$، \dots، $!D_n$، !!{formula}s نیں، تاں
    \[
    \SSubst{!A}{\subst{!D_1}{p_1}, \dots, \subst{!D_n}{p_n}} \in \Sigma,
    \]
  \item \emph{قاعدۂ اثباتِ مقدم} تحت بند اے، یعنی جے $!A$ تے $!A
      \lif !B \in \Sigma$، تاں $!B \in \Sigma$۔
  \end{enumerate}
\end{defn}

امکانی منطقاں لئی تعلقی معنیات ورتن واسطے سانوں ایہ شرط وی لاؤنی پئے گی کہ سارے امکانی نمونیاں وچ معتبر سارے !!{formula}s شامل ہون۔ پتا لگدا اے کہ ایہ شرط پوری ہو جاندی اے جے \Ax{K}\iftag{prvDiamond}{ تے~\Dual{}}{} دیاں ساریاں مثالاں
!!{derivable} ہون، تے جدوں وی !!a{formula}~$!A$، !!{derivable} ہووے، تاں~$\Box !A$ وی۔ ایہ شرطاں پوریاں کرن والی امکانی منطق نوں \emph{نارمل} آکھدے آں۔ (غیر نارمل امکانی منطقاں وی نیں، پر عام تعلقی نمونے اوہناں لئی کافی نہیں۔)

\begin{defn}
  امکانی منطق $\Sigma$ \emph{نارمل} اے جے اوہدے وچ
  \begin{align*}
    \tag{\Ax{K}} & \Box(p \lif q) \lif (\Box p \lif \Box q),
    \iftag{prvDiamond}{\\
      \tag{\Dual} & \Diamond p \liff \lnot\Box\lnot p}{}
  \end{align*}
  شامل نیں، تے اوہ \emph{لازمیت دے قاعدے} تحت بند اے، یعنی جے $!A \in
  \Sigma$، تاں $\Box !A \in \Sigma$۔
\end{defn}

دھیان دیو کہ ہمیشہ صادق لزوم اینا باریک اے کہ \emph{جہان وچ سچائی} برقرار رکھدا اے، جدکہ قاعدہ \Nec{} صرف \emph{نمونے وچ سچائی} برقرار رکھدا اے (تے ایس لئی فریم یا فریماں دی جماعت وچ منطقی اعتبار وی)۔

\begin{prop}\ollabel{prop:rk}
  ہر نارمل امکانی منطق قاعدے \RK تحت بند اے۔
  \begin{prooftree}
    \AxiomC{$!A_1 \lif (!A_2 \lif \cdots (!A_{n-1} \lif !A_n)\cdots)$}
    \RightLabel{\RK}
    \UnaryInfC{$\Box!A_1 \lif (\Box!A_2 \lif \cdots (\Box!A_{n-1}
      \lif \Box!A_n)\cdots).$}
  \end{prooftree}
\end{prop}

\begin{proof}
  استقرا نال،~$n$ اُتے۔ جے $n = 1$ تاں قاعدہ \Nec ای اے، تے ہر نارمل امکانی منطق \Nec تحت بند اے۔

  ہُن منّ لو نتیجہ $n-1$ لئی چلدا اے۔ اسیں وکھاؤنے آں کہ ایہ~$n$ لئی وی چلدا اے۔

  منّ لو
  \begin{align*}
  & !A_1 \lif (!A_2 \lif \cdots (!A_{n-1} \lif !A_n)\cdots) \in \Sigma
  \intertext{استقرائی مفروضے مطابق ساڈے کول ایہ اے}
  & \Box!A_1 \lif (\Box!A_2 \lif \cdots \Box(!A_{n-1} \lif !A_n)\cdots)
  \in \Sigma
  \intertext{کیوں جے $\Sigma$ اک نارمل امکانی منطق اے، ایس وچ~$\Ax{K}$ دیاں ساریاں مثالاں نیں، خاص طور تے ایہ}
  & \Box(!A_{n-1} \lif !A_n) \lif (\Box!A_{n-1} \lif \Box!A_n) \in \Sigma
  \intertext{قاعدۂ اثباتِ مقدم تے ہمیشہ صادق فارمولیاں توں بدل نال بنیاں مناسب مثالاں ورت کے سانوں ایہ ملدا اے}
  & \Box!A_1 \lif (\Box!A_2 \lif \cdots (\Box!A_{n-1}
  \lif \Box!A_n)\cdots) \in \Sigma.
  \end{align*}
\end{proof}

\begin{prop}\ollabel{prop:notDiamondBot}
  ہر نارمل امکانی منطق $\Sigma$ وچ~$\lnot\Diamond\lfalse$ شامل اے۔
\end{prop}

\begin{prob}
  \olref[nml][prf][nor]{prop:notDiamondBot} ثابت کرو۔
\end{prob}

\begin{prop}
  منّ لو $!A_1$، \dots، $!A_n$، !!{formula}s نیں۔ فیر اک سبھ توں چھوٹی امکانی منطق $\Sigma$ اے جیہدے وچ
  $!A_1$، \dots، $!A_n$ دیاں ساریاں مثالاں شامل نیں۔
\end{prop}

\begin{proof}
  $!A_1$، \dots، $!A_n$ دتے ہون، تاں $\Sigma$ نوں ساریاں اوہناں نارمل امکانی منطقاں دا سانجھا حصہ متعین کرو جیہناں وچ
  $!A_1$، \dots، $!A_n$ دیاں ساریاں مثالاں شامل نیں۔ سانجھا حصہ خالی نہیں، کیوں جے
  $\Frm[L]$، یعنی سارے !!{formula}s دا سیٹ، ایسی اک امکانی منطق اے۔
\end{proof}

\begin{defn}
سبھ توں چھوٹی نارمل امکانی منطق جیہدے وچ $!A_1$، \dots، $!A_n$ شامل ہون، \emph{امکانی نظام} اکھواؤندی اے تے اینوں $\Log{K} !A_1 \dots
!A_n$ لکھدے آں۔ سبھ توں چھوٹی نارمل امکانی منطق نوں~\Log{K} لکھدے آں۔
\end{defn}

\end{document}
